\begin{examplebox}
	\textbf{\textcolor{CExample}{例 1（基础）}}

	\textbf{\textcolor{CExample}{题目：}}
	已知 $\vec{a} = (3,4)$，求与 $\vec{a}$ 同方向的单位向量 $\vec{e}$。

	\textbf{\textcolor{CExample}{解题步骤：}}
	\begin{description}[style=nextline, leftmargin=2.8em, labelsep=0.5em, itemsep=0.45em]
		\item[\textbf{第一步（求向量模）：}] 计算 $|\vec{a}|$。
			\[
				\begin{aligned}
					|\vec{a}| &= \sqrt{3^2+4^2} \\
					&= \sqrt{9+16} \\
					&= \sqrt{25} \\
					&= 5.
			\end{aligned}
			\]
			\par\textbf{理由：}
			向量模的坐标公式 $(x,y)$ 的模为 $\sqrt{x^2+y^2}$。

		\item[\textbf{第二步（代入公式）：}] 代入单位向量公式 $\vec{e} = \dfrac{\vec{a}}{|\vec{a}|}$。
			\[
				\begin{aligned}
					\vec{e} &= \frac{(3,4)}{5} \\
					&= \left( \frac{3}{5}, \frac{4}{5} \right).
			\end{aligned}
			\]
			\par\textbf{理由：}
			同方向单位向量定义，各坐标除以模长。

		\item[\textbf{第三步（写出结论）：}] 得到 $\vec{e}$。
			\[
				\vec{e} = \left( \frac{3}{5}, \frac{4}{5} \right).
			\]
			\par\textbf{理由：}
			结果满足 $|\vec{e}|=1$ 且方向相同。
	\end{description}

	\textbf{\textcolor{CExample}{关键结论：}}
	\textbf{$\vec{e} = \left( \frac{3}{5}, \frac{4}{5} \right)$，长度为 1，方向与 $\vec{a}$ 相同。}

	\medskip
	\textbf{\textcolor{CExample}{例 2（稍变形）}}

	\textbf{\textcolor{CExample}{题目：}}
	已知 $A(1,2), B(4,6)$，求 $\overrightarrow{AB}$ 同方向的单位向量。

	\textbf{\textcolor{CExample}{解题步骤：}}
	\begin{description}[style=nextline, leftmargin=2.8em, labelsep=0.5em, itemsep=0.45em]
		\item[\textbf{第一步（求向量）：}] 计算 $\overrightarrow{AB}$。
			\[
				\begin{aligned}
					\overrightarrow{AB} &= (4-1,6-2) \\
					&= (3,4).
			\end{aligned}
			\]
			\par\textbf{理由：}
			终点减起点得到向量坐标。

		\item[\textbf{第二步（求模）：}] 计算 $|\overrightarrow{AB}|$。
			\[
				\begin{aligned}
					|\overrightarrow{AB}| &= \sqrt{3^2+4^2} \\
					&= 5.
			\end{aligned}
			\]
			\par\textbf{理由：}
			向量模的坐标公式 $(x,y)$ 的模为 $\sqrt{x^2+y^2}$。

		\item[\textbf{第三步（代入公式）：}] 单位向量 $\vec{e}$。
			\[
				\begin{aligned}
					\vec{e} &= \frac{\overrightarrow{AB}}{|\overrightarrow{AB}|} \\
					&= \left( \frac{3}{5}, \frac{4}{5} \right).
			\end{aligned}
			\]
			\par\textbf{理由：}
			单位向量定义。
	\end{description}

	\textbf{\textcolor{CExample}{关键结论：}}
	\textbf{$\vec{e} = \left( \frac{3}{5}, \frac{4}{5} \right)$。}

	\medskip
	\textbf{\textcolor{CExample}{例 3（综合应用）}}

	\textbf{\textcolor{CExample}{题目：}}
	已知 $\vec{a}$ 与 $\vec{b}$ 方向相同，$|\vec{a}|=3$，$|\vec{b}|=5$，$\vec{e}$ 是与它们同方向的单位向量，求 $|\vec{a}+\vec{b}|$。

	\textbf{\textcolor{CExample}{解题步骤：}}
	\begin{description}[style=nextline, leftmargin=2.8em, labelsep=0.5em, itemsep=0.45em]
		\item[\textbf{第一步（用单位向量表示）：}] 因为 $\vec{a}$ 与 $\vec{e}$ 同向，所以 $\vec{a} = |\vec{a}| \vec{e} = 3\vec{e}$；
			同理 $\vec{b} = 5\vec{e}$。
			\[
				\vec{a} = 3\vec{e}, \quad \vec{b} = 5\vec{e}.
			\]
			\par\textbf{理由：}
			若向量方向相同，则向量等于模长乘以该方向单位向量。

		\item[\textbf{第二步（求和向量）：}] 计算 $\vec{a}+\vec{b}$。
			\[
				\begin{aligned}
					\vec{a}+\vec{b} &= 3\vec{e} + 5\vec{e} \\
					&= 8\vec{e}.
			\end{aligned}
			\]
			\par\textbf{理由：}
			同类项合并。

		\item[\textbf{第三步（求模）：}] 计算 $|\vec{a}+\vec{b}|$。
			\[
				\begin{aligned}
					|\vec{a}+\vec{b}| &= |8\vec{e}| \\
					&= 8|\vec{e}| \\
					&= 8 \times 1 \\
					&= 8.
			\end{aligned}
			\]
			\par\textbf{理由：}
			模长性质 $|\lambda\vec{e}| = |\lambda||\vec{e}|$，且 $|\vec{e}|=1$。
	\end{description}

	\textbf{\textcolor{CExample}{关键结论：}}
	\textbf{$|\vec{a}+\vec{b}| = 8$。}
\end{examplebox}
